\scanpage{691}
\mspage{648}

$9^{\circ})$ \uline{Digression on group formalism}.

\textcircled{a}

Let
\[
 (1)\qquad
 \text{\uline{$G$}},\quad
 \text{\uline{$\unsure{\mathcal G}\subset G$}},\quad
 \text{\uline{$\unsure{N}\overset{\varphi_N}{\hookrightarrow}G$}},\quad
 \inserted{\text{hence}}\quad
 SN=\deleted{\unsure{N\cap\mathcal G}}\inserted{\varphi^{-1}(\unsure{\mathcal G})}\subset N
\]
be the situation of a group $G$, invariant subgroup $\unsure{\mathcal G}$, and subgroup $N$ such that
\[
 (\unsure{\alpha})\qquad
 N\xrightarrow{\text{epimorphism}}G/\unsure{\mathcal G}\deleted{\unsure{\overset{\mathrm{dfn}}{=}\gamma}},
 \qquad\text{i.e.}\quad N/SN\xrightarrow{\sim}G/\unsure{\mathcal G}
\]
\deleted{\unsure{is an epimorphism}} We are going to reconstruct this situation from the following one
\[
 (2)\qquad
 \unsure{\mathcal G},\quad N,\quad
 SN\subset N\deleted{\unsure{\operatorname{Aut}}},\quad
 N\xrightarrow{\deleted{\unsure{\varphi}}\inserted{\Theta}}\operatorname{Aut}(\unsure{\mathcal G}),
 \quad\inserted{SN\xrightarrow[\varphi=\varphi_{SN}]{\deleted{\unsure{\varphi}}}\unsure{\mathcal G}}
\]
where $\unsure{\mathcal G},N$ are groups, $SN$ a normal subgroup of $N$, $N\longrightarrow\operatorname{Aut}(\unsure{\mathcal G})$ a homomorphism of groups, i.e. an \uuline{action} of $N$ on $\unsure{\mathcal G}$; finally $SN\longrightarrow\unsure{\mathcal G}$ a homomorphism of groups. We assume that \deleted{the action}
\deleted{that $SN\longrightarrow\unsure{\mathcal G}$ commutes}
\[
 (3)\quad
 \left\{\begin{aligned}
 &\inserted{\text{a) the action}}\text{ of }SN\text{ on }\unsure{\mathcal G}\text{ \unsure{induced|induced} by that of }\unsure{\mathcal G|G},\\
 &\text{is also that deduced from }SN\longrightarrow\unsure{\mathcal G},\\
 &\text{via the adjoint operation of }\unsure{\mathcal G},\\
 &\inserted{\text{b) }SN\longrightarrow\unsure{\mathcal G}\text{ is compatible with the actions of }N.}
 \end{aligned}\right.
\]
and \deleted{\unsure{It is|I see}} \deleted{\unsure{It is|I see}} clear how (1) defines (2), let us show how a system (2) allows one to reconstruct a system (1).

\margin{690}

% Printed archive footer.
\begin{center}
690\\
\texttt{https://grothendieck.umontpellier.fr}\\
Université de Montpellier, UMR CNRS 5149, Case courrier 051, Place Eugène Bataillon\\
34095 Montpellier cedex 5 - France
\end{center}


\scanpage{692}

We thus start from \unsure{(2)}, we set
\[
 (4)\qquad S_0G=N.\unsure{\mathcal G}\qquad(\text{\unsure{Ch direct|direct product}}),
\]
we define a homomorphism
\diagram{\((5)\quad\left\{\vcenter{\hbox{\begin{tikzcd}[row sep=large,column sep=large]
 SN \arrow[r,hook,"i"] & S_0G \\
 \unsure{u} \arrow[r,mapsto] & \unsure{u}.\varphi(\unsure{u})^{-1}=\varphi(u)^{-1}.u
\end{tikzcd}}}\right.\)}

\margin{\uline{N.B.} we have \inserted{(by hypothesis (a))} for $u\in SN$, $x\in\unsure{\mathcal G}$
\[
 (*)\quad\left\{u(x)=\underset{\substack{\text{in }S_0G\\=N.\unsure{\mathcal G}}}{\underbrace{uxu^{-1}}}=\varphi(u)x\varphi(u)^{-1}\right.
\]
so if $x=\varphi(u)$, we have
\[
 u(x)=uxu^{-1}=x
\]
i.e. $x=\varphi(u)$ and $u$ commute.}

Let us check that it is a homomorphism
\deleted{\unsure{we}}
\[
 i(\unsure{u'}u)=(u'u)\varphi(u'u)^{-1}
\]
\[
\begin{aligned}
 i(u')i(u)&=u'\varphi(u')^{-1}u\varphi(u)^{-1}\\
 &\deleted{=u'\varphi(u')^{-1}u\varphi(u')}\\
 &=u'u\underset{\substack{=\varphi(u^{-1})\varphi(u'^{-1})\varphi(u)\\\text{by a)}}}{\underbrace{(u^{-1}\varphi(u'^{\unsure{-1}})u)}}\varphi(u^{-1})^{\deleted{\unsure{-1}}}\\
 &=u'u\varphi(u^{-1}u'^{-1}\deleted{uu^{-1}})=(u'u)\varphi((u'u)^{-1})=(u'u)\varphi(u'u)^{-1}\qquad\text{ok.}
\end{aligned}
\]

We have not yet used the fact that is normal in $N$, \inserted{and hypothesis (b)} this implies that the image of \unsure{(4)|(5)} is normal in $S_0G$, for \deleted{it is} we have, if $a\in N$, $u\in SN$
\[
 \deleted{\unsure{(6)(6)}}\qquad\deleted{\unsure{(*)}}\quad
 a(u.\varphi(u)^{-1})a^{-1}
 =\underbrace{aua^{-1}}_{\in SN}.\underset{=\varphi(au^{-1}a^{-1})=\varphi(au^{-1}a^{-1})^{-1}}{\underbrace{a\varphi(u^{-1})a^{-1}}}
\]
i.e. the homomorphism (5) \deleted{\unsure{is}} commutes with the actions of $N$ on $SN$ and on $S_0G\supset N$ -- thus the image of $SN$ is normalized by $N\subset S_0G$. We still have to check that it is also normalized by $\unsure{\mathcal G}$, \deleted{so let $x\in\unsure{\mathcal G}$, we have}\inserted{but we have better:}
\[
 \deleted{xu\varphi(u^{-1})x^{-1}=(xux^{-1})(x\varphi(u^{-1})x^{-1}}
\]
The image of $SN\xrightarrow{i}S_0G$ \uline{commutes} with $\unsure{\mathcal G}$, i.e.

\margin{691}

% Printed archive footer.
\begin{center}
691\\
\texttt{https://grothendieck.umontpellier.fr}\\
Université de Montpellier, UMR CNRS 5149, Case courrier 051, Place Eugène Bataillon\\
34095 Montpellier cedex 5 - France
\end{center}


\scanpage{693}
\mspage{649}

$\forall u\in SN$, $u\varphi(u)^{-1}$ induces the trivial operation on $\unsure{\mathcal G}$ --- which is indeed trivial by (3) a).

We can therefore set
\[
(6)\qquad  G\overset{\unsure{\mathrm{d\acute ef}|\mathrm{dfn}}}{=}S_0G/i(SN),
\]
we then find a commutative diagram

\noindent (7)\quad
\diagram{\begin{tikzcd}[column sep=small,row sep=small]
& N \arrow[dr,"\varphi_N"] & \\
SN \arrow[ur,hook] \arrow[dr,"\varphi_{SN}"'] & & G \\
& \unsure{\mathcal G} \arrow[ur,hook] &
\end{tikzcd}},

we see that the image of $\unsure{\mathcal G}$ in $G$ is an invariant subgroup, and that \unsure{we} have an isomorphism on the quotients:
\[
(8)\qquad  N/SN\xrightarrow{\sim}G/\unsure{\mathcal G},
\]
which implies that we also have an isomorphism
\[
(9)\qquad  \operatorname{Ker}(\varphi_{SN}:SN\longrightarrow\unsure{\mathcal G})
\xrightarrow{\sim}\operatorname{Ker}(N\xrightarrow{\varphi_N}G);
\]
\deleted{thus $\varphi_N$}

\noindent (10)\quad
\diagram{\(\begin{array}{l}
\varphi_N:N\longrightarrow G\text{ is an embedding}\\
\unsure{\Updownarrow}\\
\varphi_{SN}:SN\longrightarrow\unsure{\mathcal G}\text{ is an embedding.}
\end{array}\)}

We thus find \uline{equivalences of categories}, inverse to one another, between \inserted{\unsure{the categories}} of

\margin{692}

% Printed archive footer.
\begin{center}
692\\
\texttt{https://grothendieck.umontpellier.fr}\\
Université de Montpellier, UMR CNRS 5149, Case courrier 051, Place Eugène Bataillon\\
34095 Montpellier cedex 5 - France
\end{center}


\scanpage{694}

situations (1), and that of situations (2).

\textcircled{b}
Let us now consider a situation
\[
(11)\qquad 
\left\{\begin{array}{l}
G,\quad \unsure{\mathcal G}\subset G,\quad (\varphi_i:N_i\longrightarrow G)_{i\in I}\\[4pt]
\text{hence }SN_i=\varphi_i^{-1}(\unsure{\mathcal G})
\end{array}\right.
\]
under the condition
\[
(12)\qquad 
\forall i\in I,\quad \varphi_i:N_i\longrightarrow G/\unsure{\mathcal G}=\gamma\quad\text{epimorphism}
\]
\unsure{can} we reconstruct it from
\[
(13)\qquad 
\unsure{\mathcal G},\ (N_i,SN_i,\unsure{\varphi_i},\Theta_i)_{i\in I}
\]
where \inserted{$\forall i\in I,$} $(N_i,SN_i,\varphi_i,\Theta_i)$ satisfies the conditions (3) a) and b) --- \unsure{where|and} of course $\forall i\in I$
\[
(14)\qquad 
\deleted{\unsure{\varphi_i}}\inserted{\Theta_i}:N_i\longrightarrow\operatorname{Aut}(\unsure{\mathcal G}),\qquad
\varphi_i:SN_i\longrightarrow\unsure{\mathcal G}.
\]

It is tempting to reconstruct $G$ with the help \unsure{of|of the} $\unsure{\mathcal G}$ and of the $N_i$, as a sort of ``amalgamated sum'' relative to:

\noindent \deleted{\unsure{(15)}}\quad
\diagram{\begin{tikzcd}[column sep=small,row sep=small]
& \unsure{\mathcal G} \arrow[dr,dashed] & \\
SN_i \arrow[ur,"\varphi_i"] \arrow[dr,hook,"\mathrm{plongt}"'] & & G \\
& N_i \arrow[ur,dashed] &
\end{tikzcd}}

\margin{693}

% Printed archive footer.
\begin{center}
693\\
\texttt{https://grothendieck.umontpellier.fr}\\
Université de Montpellier, UMR CNRS 5149, Case courrier 051, Place Eugène Bataillon\\
34095 Montpellier cedex 5 - France
\end{center}


\scanpage{695}
\mspage{650}

\deleted{\uline{Proposition} \unsure{Starting|Start} from (11), hence (13)}

% small x-shaped marks follow "dire," and $\Theta_i$ in this sentence (reference marks?); not transcribed as text.
But it is more plausible to say, that for every $i$, we reconstruct \unsure{via|here} $N_i,SN_i,\varphi_i,\Theta_i$ the situation (for $i$ fixed)

\noindent \deleted{\unsure{(17)}}\inserted{\unsure{(15)}}\quad
\diagram{\begin{tikzcd}
\deleted{\unsure{\varphi_i}}N_i \arrow[r] & G_i(\unsure{\simeq}G) \\
SN_i \arrow[u,hook] \arrow[r] & \unsure{\mathcal G} \arrow[u,hook]
\end{tikzcd}}

and that we reconstruct (11) itself by giving ourselves a \uline{transitive system} \uline{of isomorphisms} \uline{between the} $G_i$ $(i\in I)$, \uline{inducing} \uline{the identity on} $\unsure{\mathcal G}$. Such a datum implies in the first place the datum of a transitive system of isomorphisms between the
\[
(16)\qquad 
\gamma_i\overset{\unsure{\mathrm{d\acute ef}|\mathrm{dfn}}}{=}N_i/SN_i
\]
--- which we \unsure{will identify} therefore with one and the same group
\[
(17)\qquad 
\begin{gathered}
\gamma,\text{ common quotient group}\\
\text{of the }N_i\text{ (by the }SN_i\text{).}
\end{gathered}
\]
Let us now identify the $G/\unsure{\mathcal G}_i=\gamma_i$ with $\gamma$
\[
(18)\qquad G_i/\unsure{\mathcal G}\simeq\gamma,
\]
and form the fiber product

\margin{694}

% Printed archive footer.
\begin{center}
694\\
\texttt{https://grothendieck.umontpellier.fr}\\
Université de Montpellier, UMR CNRS 5149, Case courrier 051, Place Eugène Bataillon\\
34095 Montpellier cedex 5 - France
\end{center}


\scanpage{696}

\[
(19)\qquad 
\unsure{g}=\underset{\text{(fiber product over }\gamma\text{)}}{\prod_{i\in I}\!{}_{\gamma}G_i}\ \supset\ \unsure{\mathcal G}^I
\]
\margin{$\unsure{g}$ is an extension of $\gamma$ \unsure{by} $\unsure{\mathcal G}^I$}

The datum of a transitive system of isomorphisms between the $G_i$, inducing the identity on the subgroups $\unsure{\mathcal G}$, and compatible with the system \deleted{\unsure{of isomorphisms|of the}} \unsure{transitive} of isomorphisms already considered between the $\gamma_i=N_i/SN_i\simeq G_i/\unsure{\mathcal G}$, is then equivalent to the datum of a subgroup $G$ of $\unsure{g}$, satisfying the conditions
\[
(20)\qquad 
\left\{\begin{array}{ll}
\text{a)}&G\cap\unsure{\mathcal G}^I=\text{diagonal }\delta(\unsure{\mathcal G})\text{ of }\unsure{\mathcal G}^I\\[5pt]
\text{b)}&\text{The projection }G\longrightarrow\gamma\text{ is an epimorphism}
\end{array}\right.
\]

In other words, we have constructed, with the help of (13) and of the transitive system of isomorphisms between the $\gamma_i$, an extension $\unsure{g}$ of their ``common value'' $\gamma$, by $\unsure{\mathcal G}^I$ --- and the additional structure, corresponding to the datum \unsure{(11)}, is the ``restriction'' of this extension to an extension of $\gamma$ by the diagonal subgroup of $\unsure{\mathcal G}^I$.

\margin{695}

% Printed archive footer.
\begin{center}
695\\
\texttt{https://grothendieck.umontpellier.fr}\\
Université de Montpellier, UMR CNRS 5149, Case courrier 051, Place Eugène Bataillon\\
34095 Montpellier cedex 5 - France
\end{center}